\chapter{Moving Around the Worksheet}
\label{ch:moving}

The worksheet is 256 columns by 65,535 rows and the screen shows seven columns
and fifty-four rows of it. This chapter is about getting from one part of it to
another.

\section{Addresses}
\index{cell!address}\index{columns!naming}

A cell is named by its column letter and its row number, in that order and with
nothing between them: \cellref{A1}, \cellref{B9}, \cellref{IV65535}.

Columns run \cellref{A} to \cellref{Z}, then \cellref{AA} to \cellref{AZ}, then
\cellref{BA} onwards, and stop at \cellref{IV} --- the 256th column.

\begin{ccfigure}
\begin{tikzpicture}[x=1pt,y=-1pt,font=\ttfamily\fontsize{8.5}{10}\selectfont]
  \foreach \i/\n in {0/A,1/B,2/C,3/{$\cdots$},4/Z,5/AA,6/AB,7/{$\cdots$},
                     8/AZ,9/BA,10/{$\cdots$},11/IU,12/IV}{
    \pgfmathsetmacro{\xx}{\i*27}
    \draw[ccink,line width=0.5pt] (\xx,0) rectangle ++(26,15);
    \node at (\xx+13,10) {\n};
  }
  \foreach \i/\n in {0/1,1/2,2/3,4/26,5/27,6/28,8/52,9/53,11/255,12/256}{
    \pgfmathsetmacro{\xx}{\i*27}
    \node[font=\sffamily\fontsize{6.8}{8}\selectfont,text=ccink] at (\xx+13,25) {\n};
  }
  \node[anchor=west,font=\sffamily\fontsize{7.4}{9}\selectfont,text=ccink]
    at (0,42) {Column letters, and the column numbers they stand for};
\end{tikzpicture}
\caption{How the columns are named}
\label{fig:colnames}
\end{ccfigure}

Rows are numbered from 1. The last cell in the worksheet is
\cellref{IV65535}.

\section{Moving one cell at a time}

\begin{cctable}{Key}{Moves the cursor}{1.55in}{3.27in}
\key{$\uparrow$} & Up one row \\
\key{$\downarrow$} & Down one row \\
\key{$\leftarrow$} & Left one column \\
\key{$\rightarrow$} & Right one column \\
\key{Return} & Down one row --- the same as \key{$\downarrow$} \\
\key{Tab} & Right one column \\
\keyc{Shift}{Tab} & Left one column \\
\end{cctable}

Movement stops at the edges. There is no wrapping from column \cellref{IV} to
column \cellref{A}, and pressing \key{$\uparrow$} on row~1 does nothing.

\section{Moving a screen at a time}

\begin{cctable}{Key}{Moves the cursor}{1.55in}{3.27in}
\key{PgUp} & Up fifty-four rows --- one screenful \\
\key{PgDn} & Down fifty-four rows \\
\end{cctable}

The distance is one screenful, whatever a screenful currently is: on a shorter
text display these keys move correspondingly less.

\section{Jumping}
\index{HOME key}\index{END key}

\begin{cctable}{Key}{Goes to}{1.55in}{3.27in}
\key{Home} & Column \cellref{A} of the row you are on \\
\keyc{Shift}{Home} & Cell \cellref{A1} \\
\key{End} & Column \cellref{IV} of the row you are on \\
\end{cctable}

\begin{cccaution}
\key{End} goes to column \cellref{IV} --- the last column that exists, not the
last column you have used. In an empty worksheet that is 255 columns of nothing.
\keyc{Shift}{Home} is the way back.
\end{cccaution}

There is no \menupath{Go to} command and no way to type an address to jump to.

\section{Scrolling}
\index{scrolling}

The screen follows the cursor; you do not scroll the screen independently.

Moving off the bottom or the top of the visible rows scrolls by one row.
\key{PgUp} and \key{PgDn} scroll by a full screen. Moving off the right-hand
edge scrolls left by whole columns, one at a time, until the cursor's column
fits on the screen completely --- so a wide column can push several narrow ones
off the left edge in a single step.

\section{Freezing the headings}
\index{freeze}\index{panes}

\menupath{Layout > Freeze} pins the rows above the cursor and the columns to its
left, so that they stay on screen while the rest scrolls.

\begin{ccprocedure}[]
\item Put the cursor on the first cell of the \emph{data} --- if the headings
are in row~1 and column~\cellref{A}, that is \cellref{B2}.
\item Choose \menupath{Layout > Freeze}.
\end{ccprocedure}

\begin{ccnote}
At \cellref{A1} there is nothing above or to the left to keep, so
\menupath{Freeze} there is taken to mean the obvious thing instead: the first
row and the first column are frozen. It is the common case, and it saves
moving the cursor to \cellref{B2} to ask for it.
\end{ccnote}

Everything above and to the left of the cursor now stays put. Choosing
\menupath{Freeze} again unfreezes: it is one command, not two, and being frozen
already is what tells it which to do.

\begin{ccnote}
\ccprod{} refuses to freeze anything that would leave nothing to scroll. Rows
are frozen only if the cursor is in the top half of the visible rows, and
columns only if at least one more column of default width would still fit
across the screen.

If neither can be frozen, the status line says \msg{Nothing to freeze here}
--- the cursor is too far down the sheet, or too far to the right. Move it up
or left and try again.
\end{ccnote}

There is no split screen, and the freeze is not saved with the workbook.

\section{Using the mouse}
\index{mouse}

If a mouse is connected, \ccprod{} shows a pointer and uses it.

\begin{cctable}{Action}{Effect}{1.55in}{3.27in}
Click in the worksheet & Moves the cell cursor to the cell you clicked. \\
Click on a worksheet tab & Switches to that worksheet. \\
Click on the menu bar & Opens the menu under the pointer. \\
Roll the wheel & Scrolls three rows per notch. Rolling towards you moves down
the worksheet. \\
Click on a row number & Nothing. There is no row selection for it to mean. \\
\end{cctable}

While you are typing into a cell the mouse is ignored entirely: clicks do not
move the cursor and do not cancel the entry. Finish or cancel the entry first.

\section{Where the cursor is}

Three things on the screen tell you, and they never disagree:

\begin{ccsteps}
\item The cell itself is drawn in reverse video.
\item Its row number and its column letter are highlighted in yellow.
\item Its address appears at the left of the formula bar and again in the
middle of the status line.
\end{ccsteps}
